%%
%% Jiaming Chu — Resume (English)
%% Template: resume.cls (Weitian LI, LPPL 1.3c; based on YAAC: Another Awesome CV)
%% Content: CC BY 4.0
%%
%% Build: latexmk resume-en.tex   (cv/.latexmkrc selects XeLaTeX, output in build/)
%% The phone number stays out of git: put \def\myphone{+86 xxx-xxxx-xxxx} in cv/private.tex
%%
\documentclass{resume}

\iconsize{\Large}
\fileinfo{\faEdit{} \today}

\name{Jiaming}{Chu}
\keywords{generative models, 3D reconstruction, world models, NeRF}

\profile{
  \privatemobile
  \email{jiaming.chu@sjtu.edu.cn}
  \github{chu-jiaming}
  \\
  \homepage{chu-jiaming.github.io}
  \university{M.Eng. student, Shanghai Jiao Tong University}
}

\begin{document}
\makeheader

%======================================================================
% Summary
%======================================================================
Working on \emph{generative models, 3D reconstruction and world models}, currently on how radiance
fields behave under imperfect camera poses and how to evaluate them faithfully.
Seeking research internships.

%======================================================================
\sectionTitle{Education}{\faGraduationCap}
%======================================================================
\begin{educations}
  \education%
    {Sep 2025}%
    {Shanghai Jiao Tong University}%
    {School of Automation and Intelligent Sensing}%
    {Electronic Information}%
    {M.Eng.}
  \separator{0.5ex}
  \education%
    {Sep 2021}%
    [Jul 2025]%
    {Harbin Institute of Technology}%
    {School of Instrumentation Science and Engineering}%
    {IntelliSense Engineering}%
    {B.Eng.\ (GPA 88.25/100)}
\end{educations}

%======================================================================
\sectionTitle{Research}{\faFlask}
%======================================================================
\begin{experiences}
  \experience%
    {2026}%
    {\textbf{PoseGuard-NeRF: When Does Joint Camera-Pose Optimization Help NeRF?} \hfill
     \link{https://chu-jiaming.github.io/poseguard-nerf/}{page} \textbullet{}
     \link{https://github.com/chu-jiaming/PoseGuard-NeRF}{code}}%
    [\begin{itemize}
      \item Controlled study of \emph{45 vanilla-NeRF runs} (9 pose-noise levels $\times$ 5 methods)
        on when joint pose optimization helps or hurts.
      \item A camera learning-rate warm-up recovers most of the loss from noisy poses:
        14.1 $\rightarrow$ \emph{22.7\,dB} PSNR at 5\% translation noise.
      \item From perfect poses it drops PSNR 27.1 $\rightarrow$ 17.9\,dB, yet a Sim(3) fit on training
        cameras recovers \emph{+5.5\,dB}: global drift, not a worse reconstruction.
      \item Released v1.0.0 with reproducible configs, GitHub Actions CI and uv environments.
    \end{itemize}]
\end{experiences}

%======================================================================
\sectionTitle{Experience}{\faBriefcase}
%======================================================================
\begin{experiences}
  \experience%
    [Jan 2023]%
    {Jul 2023}%
    {\textbf{Captain, Team Lilac 12} @ HIT Smartcar \& Innovation Club}%
    [\begin{itemize}
      \item Led vision support; built ROS packages, tuned navigation (SLAM, path planning) and wrote docs.
    \end{itemize}]
  \separator{0.5ex}
  \experience%
    [Nov 2023]%
    {Apr 2024}%
    {\textbf{Team Member} @ HIT Center of Ultra-Precision Optoelectronic Instrument Engineering}%
    [\begin{itemize}
      \item Single-beam 3-DoF ultra-precision laser interferometer: host communication and data
        visualization; the entry won a national second prize (see Awards).
    \end{itemize}]
  \separator{0.5ex}
  \experience%
    [Dec 2023]%
    {Jan 2024}%
    {\textbf{Mathematical Modeling (Team Lead)}}%
    [\begin{itemize}
      \item Built the models, wrote the code and the paper; won an ICM Honorable Mention.
    \end{itemize}]
  % [trimmed] weakly related to the target roles; uncomment if needed
  % \separator{0.5ex}
  % \experience%
  %   [Mar 2024]%
  %   {Apr 2024}%
  %   {\textbf{Visiting Student} @ Shanghai Infocrops Science \& Technology Co., Ltd.}%
  %   [\begin{itemize}
  %     \item Intelligent vision training.
  %   \end{itemize}]
\end{experiences}

%======================================================================
\sectionTitle{Other Projects}{\faCode}
%======================================================================
\begin{itemize}
  \item \link{https://github.com/chu-jiaming/qingzhou}{Qingzhou robot}
    (C/C++, Python, Jetson Nano, STM32): SLAM, navigation, visual recognition.
  \item Vision-controlled car (Python, Raspberry Pi): drives by recognizing QR codes and hand gestures.
  \item Smart-home system (Python, Raspberry Pi): real-time sensing, logging, display and RF device control.
  % [trimmed] weakly related to the target roles; uncomment if needed
  % \item Gesture-controlled car (C, STC89C51) using an MPU6050 IMU.
  % \item High-current power module PCB for precision instruments (Altium Designer, MPS).
\end{itemize}

%======================================================================
\sectionTitle{Honors \& Awards}{\faTrophy}
%======================================================================
\begin{itemize}
  \item \emph{Second National Prize}, 7th China (International) Sensor Innovation Contest \hfill 2023
  \item \emph{Provincial Second Prize}, National University Students Intelligent Car Race \hfill 2023
  \item \emph{Honorable Mention}, Interdisciplinary Contest in Modeling (ICM), COMAP \hfill 2024
\end{itemize}

%======================================================================
\sectionTitle{Skills \& Languages}{\faWrench}
%======================================================================
\begin{competences}[9.6em]
  \comptence{\icon{\faCube}Deep Learning}{PyTorch, NeRF, OpenCV, NumPy}
  \comptence{\icon{\faLaptopCode}Programming}{Python, C/C++, MATLAB, Verilog}
  \comptence{\icon{\faCog}Tools}{Linux, ROS, Git, GitHub Actions, uv, Conda, Docker, \LaTeX, SolidWorks, Altium Designer}
  \comptence{\icon{\faLanguage}Languages}{\emph{Chinese} --- native, \emph{English} --- proficient}
\end{competences}

\end{document}
